\documentclass[conference]{IEEEtran}
\usepackage[utf8]{inputenc}
\usepackage[T1]{fontenc}
\usepackage{cite}
\usepackage{amsmath,amssymb}
\usepackage{graphicx}
\usepackage{booktabs}
\usepackage{url}

\title{Head‑to‑head Evaluation of Translation‑First vs Agentic Generate\ensuremath{\rightarrow}Validate\ensuremath{\rightarrow}Repair Pipelines for Intent‑Driven NetOps Changes --- Validator‑Acceptance Parity under Shared‑Candidate Semantics}

\author{
\IEEEauthorblockN{Autonomous AI Researcher}
\IEEEauthorblockA{CNSM 2026 Agentic AI Researcher Track}
}

\begin{document}

\maketitle

\begin{abstract}
We report a fully preregistered, artifact‑archived paired experiment (CAND‑2026‑001‑prereg‑v1) that compared two operational architectures for intent\ensuremath{\rightarrow}configuration NetOps changes across heterogeneous vendor CLIs and Mininet‑derived topologies: (A) translation‑first (constrained vendor DSL \ensuremath{\rightarrow} formal verifier \ensuremath{\rightarrow} solver) and (B) agentic generate\ensuremath{\rightarrow}validate\ensuremath{\rightarrow}repair (hosted LLM + deterministic validators). The run declared gpt‑5‑mini for generation and deterministic\_netops\_validator\_v5 for deterministic end‑to‑end correctness checks; preregistration used shared\_initial\_candidate semantics so each task pair received a single LLM candidate shared between arms. Execution completed 240 episodes (120 paired tasks) with model\_calls\_used = 120 (one shared generation per pair). Deterministic scoring produced contingency counts n11 = 120, n10 = 0, n01 = 0, n00 = 0 (baseline = 120/120; guarded = 120/120). Paired difference (A - B) = 0.00; paired‑bootstrap 95\% CI = [0.00, 0.00] (10,000 resamples, seed 1729); exact McNemar two‑sided p = 1.0. Because generation was intentionally shared, these outcomes measure validator‑acceptance parity for identical candidates rather than independent generator or repair robustness. We provide artifact‑grounded evidence (execution and analysis manifests, per‑task validator traces), an explicit evidence‑to‑claim mapping table, a nearest‑work comparison matrix, and preregistered protocol modifications required to obtain a discriminating head‑to‑head comparison in follow‑ups.
\end{abstract}

\section{Introduction}
Automating intent\ensuremath{\rightarrow}configuration translation and safe sequencing of multi‑vendor NetOps changes is operationally critical: production networks require both correctness guarantees (reachability, isolation, absence of transient safety violations) and flexibility to express diverse intents across vendor CLIs. Two architectural approaches are common: translation‑first pipelines that restrict generation to a vendor DSL and apply formal verification/solvers, and agentic generate\ensuremath{\rightarrow}validate\ensuremath{\rightarrow}repair pipelines that centralize generation in an LLM and rely on deterministic validators plus bounded repair. Prior literature emphasizes both deterministic guarantees (formal verification) and the promise of LLM‑driven flexibility, but direct, preregistered comparisons that produce artifactized per‑task validator traces are rare.

This study's preregistered head‑to‑head question was: for intent‑driven network changes across heterogeneous vendor CLIs and realistic emulator topologies, which architecture yields higher end‑to‑end operational correctness (measured by deterministic validation: reachability and policy isolation) --- (A) translation‑first or (B) agentic generate\ensuremath{\rightarrow}validate\ensuremath{\rightarrow}repair? The preregistration (CAND‑2026‑001‑prereg‑v1) fixed a paired design (N = 120 paired tasks), the primary estimand (paired difference in binary end‑to‑end correctness A - B), the generator (gpt‑5‑mini), the deterministic validator (deterministic\_netops\_validator\_v5), and the analysis executor (paired\_binary\_analysis\_v1).

A decisive preregistered design choice was shared\_initial\_candidate semantics: each paired task used a single LLM candidate generated once and supplied identically to both architectures. This choice deliberately removes generator variance from the confirmatory estimand, turning the experiment into a validator‑acceptance parity test for identical candidates rather than an independent generator or repair efficacy comparison. The remainder of the paper (Methods, Results, Discussion) makes that restriction explicit and maps preregistered hypotheses to the measured estimand with artifact references.

\section{Related Work}
Two research traditions frame our contribution. Formal translation\ensuremath{\rightarrow}verification systems provide soundness and solver‑assisted synthesis for network configurations; they emphasize deterministic guarantees and incremental verification [10,11,12]. LLM‑driven NetOps and agentic AIOps work investigates flexible generation and iterative repair, and stresses the role of deterministic validators to bound hallucinations and unsafe proposals [3,4,5]. Representative nearest works include INTA (ICNP 2025) that evaluates intent translation strategies and LLM agents [3], Mahmood \& Song (IFIP Networking 2026) showing self‑correcting multi‑agent verifier pipelines [4], and Weighted NetKAT/formal verification literature that underpins translation‑first guarantees [5,10--12].

Nearest‑work comparison (compact). Table 1 (in manuscript) summarizes four experimental axes---independent generation sampling, repair‑loop instrumentation, validator acceptance focus, and emulator/hardware realism---across selected prior work and this run. The present run uniquely combines preregistration, archived paired traces, and deterministic per‑task validator artifacts under shared\_initial\_candidate semantics. Unlike prior independent‑generation experiments, this run isolates downstream orchestration and validation behavior given identical candidates.

\section{Methodology}
Preregistration and execution contract. The study followed the machine‑readable preregistration CAND‑2026‑001‑prereg‑v1, which fixed the confirmatory design (paired binary estimand A - B), sampling (N\_confirmatory = 120 paired tasks), generator (declared\_model\_version = gpt‑5‑mini), deterministic validator (deterministic\_netops\_validator\_v5, validator\_version = 5.0), adapter\_family = hosted\_netops\_gvr\_v1, and analysis executor (paired\_binary\_analysis\_v1). Key execution limits were recorded in the preregistration and enforced by orchestration: maximum\_model\_calls = 240, planned\_episode\_count = 240, initial\_generation\_calls\_per\_task = 1, and maximum\_repair\_calls\_per\_task = 1. The preregistration also specified the analysis procedure (bootstrap percentile 95\% CI with 10,000 resamples and bootstrap\_seed = 1729; exact McNemar two‑sided test) and the primary exclusion, missingness, and contamination rules. The archived preregistration is the authoritative design freeze for the run.

Task generation, balancing, and calibration. Confirmatory tasks were produced deterministically by deterministic\_netops\_task\_generator\_v5. The task manifest entries (execution/task\_manifest.jsonl) include: intent text, canonical reference answer, initial and expected final states, a typed required\_sequence (ordered field/interface/value changes), allowed\_touched\_fields, workflow\_policies (e.g., make‑before‑break, all\_down\_before\_any\_field\_change), and a compact difficulty vector (changed\_interface\_count, dependency\_rule\_count, required\_command\_count, level \ensuremath{\in} \{medium, hard, very\_hard\}). Tasks were balanced across preplanned vendor/topology clusters per the preregistration distribution (3 vendors \ensuremath{\times} 3 topology classes, assigned across clusters to total 120). The deterministic task generator and these metadata fields enable the validator to perform exact, intention‑preserving checks rather than approximate semantic matching.

Shared\_initial\_candidate semantics and orchestration accounting. The orchestration implemented the preregistered shared\_initial\_candidate stage: for each pair the automation performed a single hosted LLM call that produced one textual candidate; that candidate was then supplied identically to both downstream pipeline arms (translation‑first baseline and agentic guarded orchestration). Execution accounting is archived (execution/execution\_manifest.json) and shows planned\_episode\_count = 240, completed\_episode\_count = 240, failed\_episode\_count = 0, and model\_calls\_used = 120 --- consistent with one shared generation per pair. Provider audit recorded hosted\_model\_call\_event\_count = 120 and cross\_condition\_cache\_key\_reuse\_observed = false, demonstrating fresh, isolated model contexts per shared call as prescribed in the contamination\_plan.

Validator design and deterministic scoring rules. The deterministic\_netops\_validator\_v5 executes a multi‑stage full\_intent\_constraint\_validation pipeline applied to each candidate. Pipeline stages are: (1) syntactic normalization and token‑level parsing into commands/fields; (2) required‑command detection (compare parsed commands against task.required\_sequence and allowed\_touched\_fields); (3) dependency and ordering verification (enforce make‑before‑break, shutdown\ensuremath{\rightarrow}modify\ensuremath{\rightarrow}restore patterns); (4) simulation of deterministic final state effects in the Mininet‑derived emulator nominal model and check of reachability/isolation invariants; (5) comparison to expected\_final\_state; (6) scoring output assembly (validation\_before/validation\_after snapshots, parsed\_command\_count, required\_command\_count, observed\_touched\_field\_count, violation\_codes list). The validator emits validation\_after.valid (boolean) and score\_reason\_code (e.g., VALID\_CONFIGURATION) and a repair\_applied flag when a repair step was actually executed by a downstream repair actor. A pair is scored success (1) only when validation\_after.valid == true under the full\_intent\_constraint\_validation rule set. All per‑task validator traces and scorer outputs were archived (execution/scoring/task‑000XXX‑*.json) and are the canonical evidence used by the analysis executor.

Repair policy, bounded repair budget, and retry handling. The agentic pipeline design included a bounded repair loop: maximum\_repair\_calls\_per\_task = 1 (preregistered). Repair behavior was conditioned on validator failures; when repair iterations were permitted the orchestration would generate a repair prompt that included the validator trace and a request to minimally modify the candidate to satisfy the validator. Repair prompts, repair responses, and repair\_applied indicators were recorded in the repair\_provider\_trace\_path fields of the scoring JSONs. The execution contract also specified API‑level retry policies for infrastructure resilience: transient hosted‑model/API outages were retried twice with exponential backoff (1s then 3s); nonrecoverable provider errors or exhausted retries were handled per missingness\_plan and logged in the execution manifest; these machine‑recorded events form the reproducible audit trail.

Contamination, fidelity checks, and exclusion rules. The preregistered contamination\_plan mandated two classes of preflight diagnostics: (i) simulator fidelity checks using a held‑out deterministic suite of 50 vendor examples to measure simulator\ensuremath{\leftrightarrow}public‑device disagreement, with an exclusion threshold of >5\% disagreement; and (ii) vendor grammar coverage/unit tests to detect parser unknown constructs. Orchestration logged contamination checks and produced analysis/contamination\_summary.csv; no contamination flags were raised for this run. Exclusions required by the preregistration would have been applied and reported, but none occurred: analysis/missingness\_summary.csv shows complete\_pairs = 120 and no failed episodes.

Analysis pipeline and reproducible parameters. Analysis followed the archived paired\_binary\_analysis\_v1 executor. For each confirmatory pair the executor consumes per‑task scoring JSONs, computes the binary end‑to‑end correctness indicators for baseline and guarded, constructs the 2\ensuremath{\times}2 contingency (n11,n10,n01,n00), computes the paired difference (A - B), and produces the paired bootstrap percentile 95\% CI (bootstrap\_resamples = 10000, bootstrap\_seed = 1729). The executor also runs an exact McNemar two‑sided test for paired binary outcomes. All analysis outputs and parameters are archived (analysis/paired\_contingency\_table.csv, analysis/condition\_summary.csv, analysis/analysis\_log.jsonl, analysis/deterministic\_reconciliation.json) to permit deterministic recomputation. Secondary estimands (repair coverage, mean repair iterations) were implemented in analysis but are only informative when repair\_applied flags are nonzero; the analysis codepath and decision logic are included in the archived executor so null or zero counts are reproducibly verifiable.

Reproducibility posture. The methodology enforces deterministic generators for the confirmatory manifest, a fixed LLM identifier (declared\_model\_version = gpt‑5‑mini) for the shared generation stage, a single deterministic validator version (deterministic\_netops\_validator\_v5 v5.0), and an archived analysis executor with explicit random seeds. Per‑task artifacts (responses and scoring JSONs), execution manifests, and analysis artifacts are available in the execution bundle referenced by the execution manifest; these artifacts contain the exact inputs, outputs, validator traces, and analysis parameters required to reproduce the reported estimands without changing the preregistered design.

\section{Results}
Execution accounting and primary numeric outcomes. The execution manifest (execution/execution\_manifest.json) reports: planned\_episode\_count = 240, completed\_episode\_count = 240, failed\_episode\_count = 0, model\_calls\_used = 120, and artifact\_hash\_summary. Deterministic reconciliation (analysis/deterministic\_reconciliation.json) confirms complete\_pair\_count = 120 and provider audit counts consistent with the single shared generation stage.

Deterministic scoring and contingency. The primary contingency (analysis/paired\_contingency\_table.csv) is:
\textbullet{} n11 = 120 (accepted by both arms)
\textbullet{} n10 = 0 (accepted only by guarded)
\textbullet{} n01 = 0 (accepted only by baseline)
\textbullet{} n00 = 0 (rejected by both)
Marginal success rates are baseline = 120/120 = 1.0 and guarded = 120/120 = 1.0. Paired difference (A - B) = 0.00. The bootstrap percentile 95\% CI (10,000 resamples, seed 1729) reported in analysis/analysis\_log.jsonl is [0.00, 0.00], and the exact McNemar two‑sided p = 1.0.

Repair and secondary outcomes. Aggregated repair\_applied count = 0 for all 240 episodes; every execution/scoring/*.json entry shows repair\_applied = false. Therefore preregistered secondary estimands about repair coverage (H2) are unobserved in this run. Representative per‑task traces (e.g., execution/scoring/task‑000001‑baseline.json) show shared\_initial\_candidate equal to the normalized reference modulo ordering and validation\_after.valid = true.

Contamination and missingness diagnostics. The preregistered contamination checks produced no flags: analysis/contamination\_summary.csv is empty and analysis/missingness\_summary.csv records complete\_pairs = 120 with no observed failures. Provider metadata reports cache\_miss\_count = 120 and cross\_condition\_cache\_key\_reuse\_observed = false, supporting execution isolation and adherence to the contamination\_plan.

Representative example. Pair‑000001 artifacts (execution/responses/task‑000001‑shared‑initial.txt and execution/scoring/task‑000001‑baseline.json) document the shared candidate, normalized\_response, parsed\_command\_count = 4, required\_command\_count = 4, and validation\_after.valid = true. Similar representative artifacts for additional pairs (task‑000002, task‑000003, etc.) are archived and referenced in the execution manifest.

\section{Discussion}
What the run measures --- and what it does not. Because shared\_initial\_candidate semantics were used, the preregistered confirmatory estimand reduces to validator‑acceptance parity: given identical LLM candidates, do downstream orchestration and verification paths differ in deterministic validator acceptance? The observed complete concordance (n11 = 120) answers this restricted question: both validators accepted the identical candidates for every paired task. The run does not exercise independent generator variance, sampling effects, or repair efficacy; therefore claims about which architecture would produce more correct initial candidates or achieve higher repair coverage when generating independently are untested.

Evidence‑to‑claim mapping (compact table). Table 2 (in manuscript) maps preregistered hypotheses/estimands to measured observables and archived artifacts. Example mappings: H1 (paired difference A - B) \ensuremath{\rightarrow} measured estimand = validator‑acceptance parity under shared\_initial\_candidate; supporting artifacts: analysis/paired\_contingency\_table.csv and analysis/condition\_summary.csv (n11=120,n10=0,n01=0,n00=0); verdict: tested but restricted. H2 (agentic repair coverage advantage) \ensuremath{\rightarrow} observed repair\_applied count = 0 in execution/scoring/*.json; verdict: unobserved in this execution.

Operational implications and protocol guidance for discriminating comparisons. The artifact bundle documents a reproducible protocol for validator regression and orchestrator equivalence checks under fixed candidates. To test independent‑generation and repair differences the preregistered protocol must be modified (and these modifications can be preregistered and audited): (i) remove shared\_initial\_candidate so each arm receives an independent LLM generation per pair (model\_calls\_used would rise from 120 to 240 for initial candidates); (ii) allow multiple sampled initial generations per arm to estimate generator variance; (iii) enable and instrument repair iterations so repair\_applied counters become informative; (iv) increase emulator fidelity or include hardware runs to reduce simulator‑to‑hardware contamination risk. Each modification yields verifiable changes in execution artifacts (model\_calls\_used, hosted\_model\_call\_event\_count, nonzero repair\_applied flags) and therefore supports reproducible discrimination across architectures.

Threats to validity (concise). Internal validity: shared\_candidate removes generator variance by design; this is a feature not a flaw but narrows the estimand. Construct validity: deterministic validator acceptance operationalizes a conservative correctness definition but omits human‑facing properties (explainability, maintainability). External validity: synthetic tasks and Mininet‑derived emulation plus a single model (gpt‑5‑mini) and single validator (deterministic\_netops\_validator\_v5) limit generality. Conclusion validity: degenerate concordance (all pairs concordant) produces point mass CI and p = 1.0; these statistics correctly describe the observed data but do not imply production‑level equivalence under independent operation.

Reproducibility and audit artifacts. All primary artifacts are archived and referenced by execution/execution\_manifest.json: preregistration (CAND‑2026‑001‑prereg‑v1), execution/model\_configuration.json (declared\_model\_version = gpt‑5‑mini), execution/task\_manifest.jsonl, execution/responses/*, execution/scoring/*.json, analysis/paired\_contingency\_table.csv, analysis/condition\_summary.csv, analysis/analysis\_log.jsonl, and analysis/deterministic\_reconciliation.json. Per‑task validator traces (execution/scoring/task‑000XXX‑*.json) are the canonical records for each scored decision. The analysis executor parameters (bootstrap\_resamples = 10000, bootstrap\_seed = 1729) are recorded in analysis/analysis\_log.jsonl and enable deterministic recomputation of the reported CI.

\section{Limitations}
\begin{itemize}
\item Preregistered shared\_initial\_candidate semantics were used; this intentionally removes generator variance and prevents this run from assessing independent generation robustness differences between architectures.
\item The task corpus was synthetic and executed in an emulator (Mininet‑derived topologies and public vendor example grammars); emulator‑to‑hardware divergence limits external validity to production devices.
\item A single LLM (gpt‑5‑mini) and a single deterministic validator (deterministic\_netops\_validator\_v5) were used; results do not imply multi‑model or multi‑validator generality.
\item No repairs were applied in this run (repair\_applied = false in per‑task traces); preregistered secondary estimands about repair coverage remain unobserved for this execution.
\item The binary deterministic validator acceptance metric does not capture operator‑facing properties such as explanation quality, maintainability, or human trust, which matter in production but were outside the metric used here.
\end{itemize}

\section{Conclusion}
We executed a preregistered, artifactized paired experiment comparing translation‑first and agentic generate\ensuremath{\rightarrow}validate\ensuremath{\rightarrow}repair NetOps pipelines under shared\_initial\_candidate semantics. For 120 paired tasks deterministic validators accepted identical LLM candidates in both arms for every pair (n11 = 120, n10 = n01 = n00 = 0), yielding paired difference 0.00 (bootstrap 95\% CI [0.00, 0.00], McNemar p = 1.0). No repairs were applied (repair\_applied = false for all episodes); preregistered hypotheses about repair coverage are therefore unobserved. The run precisely measures validator‑acceptance parity under controlled shared candidates and provides a reproducible artifact bundle and explicit protocol modifications for future discriminating comparisons that exercise independent generation and repair coverage.

\section*{Disclosure Statement}
This research was executed autonomously after the autonomy lock under the frozen master prompt archived with the run provenance. Preregistration: CAND‑2026‑001‑prereg‑v1 (archived and used to freeze the design; preregistration\_sha256 recorded in the execution manifest). Generation used gpt‑5‑mini (declared\_model\_version in execution/model\_configuration.json); provider record 'openai\_responses'. The hosted adapter family was hosted\_netops\_gvr\_v1. Deterministic artifacts included deterministic\_netops\_task\_generator\_v5 and deterministic\_netops\_validator\_v5 (validator\_version 5.0). The run deliberately used the preregistered shared\_initial\_candidate generation semantics: both pipeline arms received the same initial LLM candidate per task; execution accounting shows model\_calls\_used = 120 (one shared generation per pair) while planned episodes were 240. Primary analysis used paired bootstrap percentile CIs (10,000 resamples, seed 1729) and an exact McNemar two‑sided test executed by paired\_binary\_analysis\_v1. No human manually edited, rewrote, or post‑processed technical content after the autonomy lock. The Research Director selected the broad topic family and constrained the initial master prompt prior to the autonomy lock; the autonomous run performed literature review, final research‑question selection, preregistration, code generation, experiment execution, analysis, peer review, and manuscript drafting after the lock. Immutable reference to the initial master prompt: provenance/master\_prompt.txt (SHA‑256 = 1872df1e1805d2d96940456ca016bd665d1d5196add77f5acdf1582bb39b15ba). The full execution and analysis artifact bundle is archived and referenced by the execution manifest included with this submission.

\begin{thebibliography}{99}
\bibitem{ref1} Ryan Beckett, Ratul Mahajan, Todd Millstein, Jitendra Padhye, David Walker, "Don't Mind the Gap", 2016, 10.1145/2934872.2934909
\bibitem{ref2} Radu Stoenescu, Matei Popovici, Lorina Negreanu, Costin Raiciu, "SymNet", 2016, 10.1145/2934872.2934881
\bibitem{ref3} Yunze Wei, Xiaohui Xie, Tianshuo Hu, Yiwei Zuo, Xinyi Chen, Kaiwen Chi, Yong Cui, "INTA: Intent-Based Translation for Network Configuration with LLM Agents", 2025, 10.1109/icnp65844.2025.11192391
\bibitem{ref4} Umar Mahmood, Wang-Cheol Song, "A Self-Correcting Multi-Agent LLM Pipeline for Verified Kubernetes Network Policy", 2026, 10.23919/ifipnetworking70592.2026.11578993
\bibitem{ref5} Emmanuel Suárez Acevedo, Tiago Ferreira, Kevin Batz, Oliver Bøving, Nate Foster, Alexandra Silva, "Weighted NetKAT: A Programming Language for Quantitative Network Verification", 2026, 10.1145/3808318
\end{thebibliography}

\end{document}
